% Filename: lista8.tex
% 
% This code is part of 'Solutions for MS550, M\'{e}todos de Matem\'{a}tica Aplicada I, and F520, M\'{e}todos Matem\'{a}ticos da F\'{i}sica I'
% 
% Description: This file corresponds to the solutions of homework sheet 8.
% 
% Created: 21.05.12 08:46:45 PM
% Last Change: 14.07.12 10:19:05 AM
% 
% Authors:
% - Raniere Silva (2012): initial version
% 
% Copyright (c) 2012 Raniere Silva <r.gaia.cs@gmail.com>
% 
% This work is licensed under the Creative Commons Attribution-ShareAlike 3.0 Unported License. To view a copy of this license, visit http://creativecommons.org/licenses/by-sa/3.0/ or send a letter to Creative Commons, 444 Castro Street, Suite 900, Mountain View, California, 94041, USA.
%
% This work is distributed in the hope that it will be useful, but WITHOUT ANY WARRANTY; without even the implied warranty of MERCHANTABILITY or FITNESS FOR A PARTICULAR PURPOSE.
%
\documentclass[a4paper,12pt, leqno, answers]{exam}
% Customiza\c{c}\~{a}o da classe exam
\newcommand{\mycheader}{Lista 8 - Problema de Sturm-Liouville}
\header{MS550, F520}{\mycheader}{\thepage/\numpages}
\headrule
\footer{Dispon\'{i}vel em \\% Filename: repository.tex
% 
% This code is part of 'Solutions for MS550, M\'{e}todos de Matem\'{a}tica Aplicada I, and F520, M\'{e}todos Matem\'{a}ticos da F\'{i}sica I'
% 
% Description: This file keeps the url of the repository.
% 
% Created: 07.03.12 04:00:00 PM
% Last Change: 30.05.12 04:40:25 PM
% 
% Authors:
% - Raniere Silva (2012): initial version
% 
% Copyright (c) 2012 Raniere Silva <r.gaia.cs@gmail.com>
% 
% This work is licensed under the Creative Commons Attribution-ShareAlike 3.0 Unported License. To view a copy of this license, visit http://creativecommons.org/licenses/by-sa/3.0/ or send a letter to Creative Commons, 444 Castro Street, Suite 900, Mountain View, California, 94041, USA.
%
% This work is distributed in the hope that it will be useful, but WITHOUT ANY WARRANTY; without even the implied warranty of MERCHANTABILITY or FITNESS FOR A PARTICULAR PURPOSE.
%
\url{https://github.com/r-gaia-cs/solucoes_listas_metodos}
}{}{Reportar erros para \\\href{mailto:r.gaia.cs@gmail.com}{r.gaia.cs@gmail.com}
}
\footrule 
\pagestyle{headandfoot}
\renewcommand{\solutiontitle}{\noindent\textbf{Solu\c{c}\~{a}o:}\enspace}
\SolutionEmphasis{\itshape}
\unframedsolutions
\pointname{}

% Filename: paper_size.tex
%
% This code is part of 'Solutions for MS550, M\'{e}todos de Matem\'{a}tica Aplicada I, and F520, M\'{e}todos Matem\'{a}ticos da F\'{i}sica I'
% 
% Description: This file corresponds to the paper size output.
% 
% Created: 07.03.12 04:00:00 PM
% Last Change: 30.05.12 04:40:25 PM
% 
% Authors:
% - Raniere Silva (2012): initial version
% 
% Copyright (c) 2012 Raniere Silva <r.gaia.cs@gmail.com>
% 
% This work is licensed under the Creative Commons Attribution-ShareAlike 3.0 Unported License. To view a copy of this license, visit http://creativecommons.org/licenses/by-sa/3.0/ or send a letter to Creative Commons, 444 Castro Street, Suite 900, Mountain View, California, 94041, USA.
%
% This work is distributed in the hope that it will be useful, but WITHOUT ANY WARRANTY; without even the implied warranty of MERCHANTABILITY or FITNESS FOR A PARTICULAR PURPOSE.
%
% Para impress\~{a}o
\usepackage[top=3cm, bottom=3cm, left=2cm, right=2cm]{geometry}

% Para ereaders (Kindle, Nook, Kobo, ...) and tablets (iPad, GalaxyTab, ...)
% \usepackage[papersize={180mm,240mm},margin=2mm]{geometry}
% \sloppy


% Pacotes
\usepackage[utf8]{inputenc}
\usepackage[T1]{fontenc}
\usepackage[brazil]{babel}
\usepackage{amsmath}
\usepackage{amsfonts}
\usepackage{amssymb}
\usepackage{hyperref}
\usepackage{graphicx}

% Customiza\c{c}\~{a}o do pacote amsmath
\allowdisplaybreaks[4]

% Novos ambientes
% \newenvironment{fwsolution}{\begin{EnvFullwidth}\begin{TheSolution}}{\end{TheSolution}\end{EnvFullwidth}}

% Novos comandos
\newcommand{\devp}[2]{\frac{\partial #1}{\partial #2}}
\newcommand{\grad}{\mbox{grad }}
\newcommand{\diver}{\mbox{div }}
\newcommand{\rot}{\mbox{rot }}

\begin{document}
%cover
\thispagestyle{empty}
% Filename: cover.tex
% This code is part of 'Solutions for MS550, M\'{e}todos de Matem\'{a}tica Aplicada I, and F520, M\'{e}todos Matem\'{a}ticos da F\'{i}sica I'
% 
% Description: This file corresponds to the cover.
% 
% Created: 30.05.12 04:40:25 PM
% Last Change: 31.05.12 10:11:55 PM
% 
% Authors:
% - Raniere Silva (2012): initial version
% 
% Copyright (c) 2012 Raniere Silva <r.gaia.cs@gmail.com>
% 
% This work is licensed under the Creative Commons Attribution-ShareAlike 3.0 Unported License. To view a copy of this license, visit http://creativecommons.org/licenses/by-sa/3.0/ or send a letter to Creative Commons, 444 Castro Street, Suite 900, Mountain View, California, 94041, USA.
%
% This work is distributed in the hope that it will be useful, but WITHOUT ANY WARRANTY; without even the implied warranty of MERCHANTABILITY or FITNESS FOR A PARTICULAR PURPOSE.
%
\begin{center}
    \LARGE{Solu\c{c}\~{o}es para MS550, M\'{e}todos de Matem\'{a}tica Aplicada I, e F520, M\'{e}todos Matem\'{a}ticos da F\'{i}sica I}
    
    \Large{\mycheader}
\end{center}
\vspace{.5\textheight}

\begin{tabular}{|p{.9\textwidth}|}
\hline
Este trabalho foi licenciado com a Licen\c{c}a Creative Commons Atribui\c{c}\~{a}o - CompartilhaIgual 3.0 N\~{a}o Adaptada. Para ver uma c\'{o}pia desta licen\c{c}a, visite \url{http://creativecommons.org/licenses/by-sa/3.0/} ou envie um pedido por carta para Creative Commons, 444 Castro Street, Suite 900, Mountain View, California, 94041, USA.
\begin{center}
\includegraphics[scale=1]{cc-by-sa.png}
\end{center}
Este trabalho \'{e} distribuido na esperança que possa ser \'{u}til, mas SEM NENHUMA GARANTIA; sem uma garantia implicita de ADEQUA\c{C}\~{A}O a qualquer MERCADO ou APLICA\c{C}\~{A}O EM PARTICULAR.
\\ \hline
\end{tabular}

\newpage
\setcounter{page}{1}
\nocite{Riley:2006:Mathematical}Uma equa\c{c}\~{a}o diferencial ordin\'{a}ria linear de ordem $n$ cuja vari\'{a}vel independente \'{e} $x$ e a vari\'{a}vel dependente \'{e} $y$ possue a forma
\begin{align*}
    a_n(x) \frac{\mathrm{d}^ny}{\mathrm{d}x^n} + a_{n - 1}(x) \frac{\mathrm{d}^{n - 1}y}{\mathrm{d}x^{n - 1}} + \ldots + a_1(x) \frac{\mathrm{d}y}{\mathrm{d}x} + a_0(x) y &= f(x). \tag{EDOL} \label{eq:EDOL}
\end{align*}
A solu\c{c}\~{a}o geral de \eqref{eq:EDOL} dever\'{a} conter $n$ constantes que podem ser determinadas se $n$ condi\c{c}\~{o}es de contorno forem fornecidas.

Para resolver qualquer \eqref{eq:EDOL} precisamos, primeiramente, encontrar a solu\c{c}\~{a}o geral para equa\c{c}\~{a}o homog\^{e}nea associada, i.e., $f(x) = 0$. Para determinar a solu\c{c}\~{a}o geral da equa\c{c}\~{a}o homog\^{e}nea associada precisamos encontrar $n$ fun\c{c}\~{o}es linearmente independentes que satisfa\c{c}am a equa\c{c}\~{a}o. Uma vez que encontramos as fun\c{c}\~{o}es, a solu\c{c}\~{a}o geral da equa\c{c}\~{a}o homog\^{e}nea associada \'{e} dada pela combina\c{c}\~{a}o linear dessas fun\c{c}\~{o}es, i.e., se as $n$ fun\c{c}\~{o}es s\~{a}o $y_1(x), y_2(x), \ldots, y_n(x)$ ent\~{a}o a solu\c{c}\~{a}o geral da equa\c{c}\~{a}o homog\^{e}nea associada \'{e}
\begin{align*}
    y_h(x) &= c_1 y_1(x) + c_2 y_2(x) + \ldots + c_n y_n(x),
\end{align*}
onde $c_i$, $i = 1, \ldots, n$, s\~{a}o constantes arbitr\'{a}rias que podem ser determinadas pelas condi\c{c}\~{o}es de contorno.

Se $f(x) = 0$ em \eqref{eq:EDOL} ent\~{a}o a solu\c{c}\~{a}o geral j\'{a} foi encontrada. Caso contr\'{a}rio, se $f(x) \neq 0$ ent\~{a}o a solu\c{c}\~{a}o geral \'{e} dado por
\begin{align*}
    y(x) &= y_h(x) + y_p(x),
\end{align*}
onde $y_p(x)$ \'{e} uma solu\c{c}\~{a}o particular qualquer da \eqref{eq:EDOL} e linearmente independente de $y_h(x)$.

Se \eqref{eq:EDOL} tiver coeficientes constantes, i.e., se $a_i(x)$ for constante ao inv\'{e}s de fun\c{c}\~{a}o de $x$, ent\~{a}o temos
\begin{align*}
    a_n \frac{\mathrm{d}^ny}{\mathrm{d}x^n} + a_{n - 1} \frac{\mathrm{d}^{n - 1}y}{\mathrm{d}x^{n - 1}} + \ldots + a_1 \frac{\mathrm{d}y}{\mathrm{d}x} + a_0 y &= f(x). \tag{EDOLCC} \label{eq:EDOLCC}
\end{align*}
O m\'{e}todo padr\~{a}o para encontrar $y_h(x)$ \'{e} tentar solu\c{c}\~{o}es da forma $y(x) = A \exp(r x)$, $A \neq 0$. Substituindo $y(x) = A \exp(r x)$ em \eqref{eq:EDOLCC} e notando que $\exp(r x) \neq 0$ obtemos uma equa\c{c}\~{a}o polinomial em $r$ de ordem $n$ representada por
\begin{align*}
    a_n r^n + a_{n - 1} \lambda^{n - 1} + \ldots + a_1 \lambda_1 + a_0 = 0
\end{align*}
que \'{e} denominada equa\c{c}\~{a}o indicial. Ao resolver a equa\c{c}\~{a}o indicial ocorre um dos tr\^{e}s casos abaixo:
\begin{description}
    \item[Todas as ra\'{i}zes s\~{a}o reais e distintas.] Neste caso as $n$ solu\c{c}\~{o}es para \eqref{eq:EDOLCC} s\~{a}o da forma $\exp(r_i x)$, $i = 1, \ldots, n$.
    \item[Alguma das ra\'{i}zes \'{e} complexa.] Se todos os coeficientes em \eqref{eq:EDOLCC}, $a_i$, forem reais ent\~{a}o as ra\'{i}zes complexas ocorrem aos pares conjugados, i.e., se $\alpha + i \beta$ \'{e} uma raiz da equa\c{c}\~{a}o indicial ent\~{a}o $\alpha - i \beta$ tamb\'{e}m \'{e} e podemos escrever
        \begin{align*}
            c_1 \exp\left( (\alpha + i \beta) x \right) + c_2 \exp\left( (\alpha - i \beta) x \right) &= \exp(\alpha x) \left( d_1 \cos(\beta x) + d_2 \sin(\beta x) \right) \\
            &= A \exp(\alpha x) \sin(\beta x + \phi) \\
            &\quad \text{ ou } A \exp(\alpha x) \cos(\beta x + \phi)
        \end{align*}
    \item[Existem ra\'{i}zes repetidas.] Se, por exemplo, $\lambda_1$ for $k$, $k > 1$, solu\c{c}\~{a}o da equa\c{c}\~{a}o indicial ent\~{a}o n\~{a}o encontramos $n$ solu\c{c}\~{o}es linearmente independentes. Mas por substitui\c{c}\~{a}o direta em \eqref{eq:EDOLCC} verifica-se que
        \begin{align*}
            x \exp(\lambda_1 x), x^2 \exp(\lambda_1 x), \ldots, x^{k - 1} \exp(\lambda_1 x)
        \end{align*}
        s\~{a}o solu\c{c}\~{o}es e pelo Wronskiano que s\~{a}o linearmente independentes, logo temos nosso conjunto de solu\c{c}\~{o}es linearmente independentes. Este racioc\'{i}nio \'{e} v\'{a}lido para qualquer uma das ra\'{i}zes repetidas.
\end{description}

Se \eqref{eq:EDOL} tiver coeficientes lineares, i.e., se $a_i(x)$ for uma fun\c{c}\~{a}o linear, ent\~{a}o temos
\begin{align*}
    a_n x^n \frac{\mathrm{d}^ny}{\mathrm{d}x^n} + a_{n - 1} x^{n - 1} \frac{\mathrm{d}^{n - 1}y}{\mathrm{d}x^{n - 1}} + \ldots + a_1 x \frac{\mathrm{d}y}{\mathrm{d}x} + a_0 y &= f x. \tag{Euler} \label{eq:Euler}
\end{align*}
que pode ser resolvida de maneira similar a \eqref{eq:EDOLCC} ao utilizar a substitui\c{c}\~{a}o $x = \exp(t)$. Se $f(x) = 0$, ent\~{a}o podemos tentar solu\c{c}\~{o}es da forma $y(x) = x^r$.

Para as demais \eqref{eq:EDOL} n\~{a}o existe um m\'{e}todo geral. Algumas t\'{e}cnicas poss\'{i}veis s\~{a}o: solu\c{c}\~{o}es exatas, varia\c{c}\~{a}o de par\^{a}metro, \ldots Tamb\'{e}m \'{e} poss\'{i}vel procurar por solu\c{c}\~{o}es na forma de uma s\'{e}rie.

Algumas equa\c{c}\~{o}es diferenciais ordin\'{a}rias lineares de segunda ordem s\~{a}o importantes:
\begin{align}
    \begin{split}
        & \text{Hipergeom\'{e}trica:} \\
        & \qquad z (1 - z) y'' + \left[ \nu - \left( a + b + 1 \right) z \right] y' - a b \nu = 0, \\
        & \qquad y(x) = A \,_2F_1(a, b, c; x) + B x^{1 - x} \,_2F_1(a - c + 1, b - c + 1, 2 - c; x), % c, a - b, c - a - b \notin \mathbb{Z}, | x | < 1,
    \end{split} \label{eq:hiperg} \\
    \begin{split}
        & \text{Hipergeom\'{e}trica confluente:} \\
        & \qquad z y'' + (c - z) y' - a y = 0,
    \end{split} \label{eq:hiperg_conf} \\
    \begin{split}
        & \text{Legendre:} \\
        & \qquad \left( 1 - z^2 \right) u'' - 2 z y' + l \left( l + 1 \right) y = 0, \\
        & \qquad y(x) = c_1 P_l(x) + c_2 Q_l(x), l \in \mathbb{Z},
    \end{split} \label{eq:legendre} \\
    \begin{split}
        & \text{Legendre Associado:} \\
        & \qquad \left( 1 - z^2 \right) y'' - 2 z y' + \left[ l \left( l + 1 \right) - m^2 / \left( 1 - z^2 \right) \right] y = 0, \\
        & \qquad y(x) = c_1 P_l^m(x) + c_2 Q_l^m(x), l, m \in \mathbb{Z},
    \end{split} \label{eq:legendre_ass} \\
%     \begin{split}
%         & \text{Chebyshev:} \\
%         & \qquad (1 - z^2) y'' - z y ' + \nu^2 y = 0,
%     \end{split} \label{eq:chenyshev} \\
    \begin{split}
        & \text{Bessel:} \\
        & \qquad z^2 y'' + z y' + \left( z^2 - \nu^2 \right) y = 0, \\
        & \qquad y(x) = c_1 J_\nu(x) + c_2 J_{-\nu}(x), \nu \notin \mathbb{Z}, \\
        & \qquad y(x) = c_1 J_\nu(x) + c_2 Y_\nu(x), \nu \in \mathbb{Z},
    \end{split} \label{eq:bessel} % \\
%     \begin{split}
%         & \text{Laguerre:} \\
%         & \qquad z y'' + \left( 1 - z \right) y' + \nu y = 0,
%     \end{split} \label{eq:laguerre} \\
%     \begin{split}
%         & \text{Laguerre associado:} \\
%         & \qquad z y'' + \left( m + 1 - z \right) y' + \left( \nu - m \right) y = 0,
%     \end{split} \label{eq:laguerre_ass} \\
%     \begin{split}
%         &\text{Hermite:} \\
%         & \qquad y'' - 2 z y' + 2 \nu y = 0,
%     \end{split} \label{eq:hermite} \\
%     \begin{split}
%         & \text{Oscilador harm\^{o}nico simples:} \\
%         & \qquad y'' + \omega^2 y = 0.
%     \end{split} \label{eq:harm_osc}
\end{align}
\begin{questions}
    \question Encontre os autovalores e autofun\c{c}\~{o}es dos seguintes problemas de Sturm-Liouville:
    \begin{parts}
        \part $\begin{cases}
            y'' + \lambda y = 0, & 0 < x < 1, \\
            y(0) + y'(0) = 0, \\
            y(1) = 0.
        \end{cases}$
        \begin{solution}
            Para $\lambda = 0$ temos que
            \begin{align*}
                y'' + \lambda y = y'' = 0 \Rightarrow y = a x + b.
            \end{align*}
            Quando $x = 1$ temos
            \begin{align*}
                0 &= y(1) = a + b \Rightarrow a = -b.
            \end{align*}
            Quando $x = 0$ temos
            \begin{align*}
                0 &= y(0) + y'(0) = -a + a = 0.
            \end{align*}
            Portanto, para $\lambda = 0$ temos que $y(x) = a x - a = a (x - 1)$ \'{e} solu\c{c}\~{a}o.

            Para $\lambda > 0$, fazemos $\lambda = k^2$ com $k > 0$, e assim temos $y'' + k^2 y = 0$. Logo $y = \exp(rx)$ \'{e} solu\c{c}\~{a}o e temos
            \begin{align*}
                r^2 + k^2 = 0 \Rightarrow r = \pm i k.
            \end{align*}
            Ent\~{a}o
            \begin{align*}
                y(x) &= A \sin(k x) + B \cos(k x), \\
                y'(x) &= A k \cos(k x) - B k \sin(k x), \\
                y(x) + y'(x) &= (A - B k) \sin(k x) + (A k + B) \cos(k x).
            \end{align*}
            Quando $x = 0$ temos
            \begin{align*}
                0 &= y(0) + y'(0) = A k + B \Rightarrow A = - B k^{-1}, k \neq 0.
            \end{align*}
            Quando $x = 1$ temos
            \begin{align*}
                0 &= y(1) = A \sin(k) + B \cos(k) = -B k^{-1} \sin(k) + B \cos(k) \\
                &= B k^{-1} \left( - \sin(k) + k \cos(k) \right) \Rightarrow k = \tan(k) \text{ ou } B = 0.
            \end{align*}
            Portanto,
            \begin{align*}
                y(x) &= - (\tan(k))^{-1} \sin(\tan(k) x) + \cos(\tan(k) x), k = 1, 2, 3, \ldots.
            \end{align*}

            Para $\lambda < 0$, fazemos $\lambda = - k^2$ com $k > 0$ e assim temos $y'' - k^2 y = 0$. Logo, $y = \exp(rx)$ \'{e} solu\c{c}\~{a}o e temos
            \begin{align*}
                r^2 - k^2 = 0 \Rightarrow r = \pm k.
            \end{align*}
            Ent\~{a}o
            \begin{align*}
                y(x) &= C \cosh(k x) + D \sinh(k x).
            \end{align*}
            Quando $x = 1$ temos
            \begin{align*}
                0 &= y(1) = C \cosh(k) + D \sinh(k) \Rightarrow C = .
            \end{align*}
            Quando $x = 0$ temos
            \begin{align*}
                0 & = y(0) + y'(0) = C k \sinh(0) + D k \cosh(0) + C \cosh(0) + D \sinh(0) \\
                &= D k + C \Rightarrow C = - D.
            \end{align*}
            Portanto, existe apenas a solu\c{c}\~{a}o trivial.
        \end{solution}

        \part $\begin{cases}
            x^2 y'' + x y' + \lambda y = 0, & 1 < x < e, \\
            y(1) = 0, \\
            y(e) = 0.
        \end{cases}$
        \begin{solution}
            Para $\lambda = 0$, temos que
            \begin{align*}
                x^2 y'' + x y' + \lambda y = x^2 y'' + x y' = 0.
            \end{align*}
            Logo, $y = x^r$ \'{e} solu\c{c}\~{a}o,
            \begin{align*}
                0 &= r (r - 1) + r = r^2 - r + r = r^2 \Rightarrow r = 0
            \end{align*}
            e portanto
            \begin{align*}
                y(x) = A x^0 + B \ln(x) x^0 = A + B \ln(x).
            \end{align*}
            Quando $x = 1$ temos
            \begin{align*}
                0 &= y(1) = A + B \ln(1) \Rightarrow A = 0.
            \end{align*}
            Quando $x = \exp(1)$ temos
            \begin{align*}
                0 & y(\exp(1)) = B \ln(\exp(1)) \Rightarrow B = 0.
            \end{align*}
            Portanto, para $\lambda = 0$ temos apenas a solu\c{c}\~{a}o trivial.

            Para $\lambda > 0$, fazendo $\lambda = k^2$ com $k > 0$, temos que
            \begin{align*}
                x^2 y'' + x y' + \lambda y = x^2 y'' + x y' + k^2 y = 0.
            \end{align*}
            Logo, $y = x^r$ \'{e} solu\c{c}\~{a}o,
            \begin{align*}
                0 &= r (r - 1) + r + k^2 = r^2 + k^2 \Rightarrow r = \pm k i
            \end{align*}
            e portanto
            \begin{align*}
                y(x) = A x^{-k i} + B x^{k i}.
            \end{align*}
            Quando $x = 1$ temos
            \begin{align*}
                0 &= y(1) = A + B \Rightarrow A = -B.
            \end{align*}
            Quando $x = \exp(1)$ temos
            \begin{align*}
                0 &= y(\exp(1)) = A (\exp(-k i) - \exp(k i)) = - A (\exp(k i) - \exp(-k i)) \\
                &= -2 A \sin(k) \Rightarrow k = n \pi, n = 1, 2, \ldots \text{ ou } A = 0.
            \end{align*}
            Portanto, para $\lambda > 0$ temos al\'{e}m da solu\c{c}\~{a}o trivial
            \begin{align*}
                y_n(x) &= x^{- n \pi i} - x^{n \pi i} \\
                &= \exp(\ln(x^{- n \pi i})) - \exp(\ln(x^{n \pi i})) \\
                &= \cos(n \pi \ln(x)) + \sin(n \pi \ln(x)) \\
                &= \sin(n \pi \ln(x)).
            \end{align*}

            Para $\lambda < 0$, fazendo $\lambda = - k^2$ com $k > 0$, temos que
            \begin{align*}
                x^2 y'' + x y' + \lambda y = x^2 y'' + x y' - k^2 y = 0.
            \end{align*}
            Logo, $y = x^r$ \'{e} solu\c{c}\~{a}o,
            \begin{align*}
                0 &= r (r - 1) + r - k^2 = r^2 - k^2 \Rightarrow r = \pm k
            \end{align*}
            e portanto
            \begin{align*}
                y(x) = A x^{k} + B x^{-k}.
            \end{align*}
            Quando $x = 1$ temos
            \begin{align*}
                0 &= y(1) = A + B \Rightarrow A = - B.
            \end{align*}
            Quando $x = \exp(1)$ temos
            \begin{align*}
                0 &= y(\exp(1)) = A \exp(k) + B \exp(-k) \Rightarrow \text{n\~{a}o tem solu\c{c}\~{a}o.}
            \end{align*}
            Portanto, para $\lambda < 0$ n\~{a}o temos solu\c{c}\~{a}o.

            Assim sendo, temos, al\'{e}m da solu\c{c}\~{a}o tivial, $y_n = \sin(n \pi \ln(x))$ quando $\lambda > 0$.
        \end{solution}

        \part[P3 de 2006] $\begin{cases}
            y'' + y' + (1 + \lambda) y = 0, & 0 < x < 1, \\
            y(0) = 0, \\
            y(1) = 0.
        \end{cases}$
        \begin{solution}
            Para $y = \exp(r x)$ temos que
            \begin{align*}
                r^2 - r + \left( 1 + \lambda \right) &= 0
            \end{align*}
            e portanto $r = (-1/2) \pm \sqrt{-3/4 - \lambda}$. Logo as solu\c{c}\~{o}es n\~{a}o triviais ocorrem quando $-3/4 - \lambda < 0$.

            Para $-3/4 - \lambda = - k^2$, $k > 0$, temos
            \begin{align*}
                r &= -1/2 \pm \sqrt{-k^2} \\
                &= -1/2 \pm i k.
            \end{align*}
            Logo,
            \begin{align*}
                y(x) &= A_1 \exp\left( (-1/2 + ik) x \right) + A_2 \exp\left( (-1/2 - i k) x \right) \\
                &= \exp(-x/2) \left( A_1 \exp(i k x) + A_2 \exp(-i k x) \right) \\
                &= \exp(-x/2) \left( B_1 \cos(k x) + B_2 \sin(k x) \right).
            \end{align*}
            Para $x = 0$,
            \begin{align*}
                0 = y(0) = \exp(0) B_1 \Rightarrow B_1 = 0,
            \end{align*}
            e para $x = 1$,
            \begin{align*}
                0 = y(1) = \exp(1) B_2 \sin(k) \Rightarrow \sin(k) = 0 \Rightarrow k = n \pi, n = 1, 2, \ldots.
            \end{align*}
            Portanto,
            \begin{align*}
                -3/4 - \lambda = -(n \pi)^2 \Rightarrow \lambda_n = n^2 \pi^2 - 3/4, n = 1, 2, \ldots.
            \end{align*}
            Se $B_1 = 0$, ent\~{a}o $k = n \pi$ e portanto
            \begin{align*}
                y_n(x) = \exp(-x/2) \sin(n \pi x), n = 1, 2, \ldots.
            \end{align*}
        \end{solution}

        \part $\begin{cases}
            y'' + \lambda y = 0, & -1 < x < 1, \\
            y(-1) = y(1), \\
            y'(-1) = y'(1).
        \end{cases}$
        \begin{solution}
            Para $\lambda = 0$ temos que
            \begin{align*}
                y'' + \lambda y = y'' = 0.
            \end{align*}
            Logo,
            \begin{align*}
               y(x) &= A x + B, \\
               y'(x) &= A.
            \end{align*}
            Quando $x = -1$ temos
            \begin{align*}
                y(-1) &= -A + B, \\
                y'(-1) &= A.
            \end{align*}
            Quando $x = 1$ temos
            \begin{align*}
                y(1) &= A + B, \\
                y'(1) &= A.
            \end{align*}
            Portanto, para $\lambda = 0$ temos apenas a solu\c{c}\~{a}o tivial.

            Para $\lambda > 0$, fazendo $\lambda = k^2$ com $k > 0$,  temos que
            \begin{align*}
                y'' + \lambda y = y'' + k^2 y = 0.
            \end{align*}
            Logo,
            \begin{align*}
                y(x) &= A \cos(k x) + B \sin(k x), \\
                y'(x) &= - A k \sin(k x) + B k \cos(k x).
            \end{align*}
            Quando $x = -1$ temos
            \begin{align*}
                y(-1) &= A \cos(-k) + B \sin(-k) = A \cos(k) - B \sin(k), \\
                y'(-1) &= -A k \sin(-k) + B k \cos(-k) = A k \sin(k) + B k \cos(k).
            \end{align*}
            Quando $x = 1$ temos
            \begin{align*}
                y(1) &= A \cos(k) + B \sin(k), \\
                y'(1) &= - A k \sin(k) + B k \cos(k).
            \end{align*}
            Portanto, para $\lambda > 0$ temos, al\'{e}m da solu\c{c}\~{a}o trivial, $y_n(x) = \cos(n \pi x)$ e $y_n(x) = \sin(n \pi x)$ onde $\lambda n = n^2 \pi^2$ e $n = 1, 2, \ldots$.

            Para $\lambda < 0$, fazendo $\lambda = - k^2$ com $k > 0$, temos que
            \begin{align*}
                y'' + \lambda y = y'' - k^2 y = 0.
            \end{align*}
            Logo,
            \begin{align*}
                y(x) &= A \cosh(k x) + B \sinh(k x), \\
                y'(x) &= - A k \sinh(k x) + B k \cosh(k x).
            \end{align*}
            Quando $x = -1$ temos
            \begin{align*}
                y(-1) &= A \cosh(-k) + B \sinh(-k), \\
                y'(-1) &= -A k \sinh(-k) + B k \cosh(-k).
            \end{align*}
            Quando $x = 1$ temos
            \begin{align*}
                y(1) &= A \cosh(k) + B \sinh(k), \\
                y'(1) &= - A k \sinh(k) + B k \cosh(k).
            \end{align*}
            Portanto, para $\lambda < 0$ n\~{a}o temos solu\c{c}\~{a}o.

            Assim sendo, temos, al\'{e}m da solu\c{c}\~{a}o tivial, $y_n(x) = \cos(n \pi x)$ e $y_n(x) = \sin(n \pi x)$ onde $\lambda n = n^2 \pi^2$ e $n = 1, 2, \ldots$ quando $\lambda > 0$.
        \end{solution}

        \part $\begin{cases}
            y'' - 3y' + 3(1 + \lambda)y = 0, & 0 < x < \pi, \\
            y(0) = 0, \\
            y(\pi) = 0.
        \end{cases}$
        \begin{solution}
            Sabendo que $y = \exp(r x)$ \'{e} solu\c{c}\~{a}o temos que
            \begin{align*}
                0 &= r^2 - 3 r + 3 (1 + \lambda) \Rightarrow r = \left( 3 \pm \sqrt{-3 - 12 \lambda} \right) 2^{-1}.
            \end{align*}

            Para $-3 - 12 \lambda = 0$, temos que
            \begin{align*}
                y(x) = A \exp(3 x 2^{-1}) + B x \exp(3 x 2^{-1}).
            \end{align*}
            Quando $x = 0$ temos
            \begin{align*}
                0 &= y(0) = A.
            \end{align*}
            Quando $x = \pi$ temos
            \begin{align*}
                0 &= y(\pi) = B \pi \exp(3 \pi 2^{-1}) \Rightarrow B = 0.
            \end{align*}
            Portanto, para $-3 - 12 \lambda = 0$ temos apenas a solu\c{c}\~{a}o trivial.

            Para $-3 - 12 \lambda > 0$, temos que
            \begin{align*}
                y(x) &= C \exp\left( (3 + k) x 2^{-1} \right) + D \exp\left( (3 - k) x 2^{-1} \right).
            \end{align*}
            Quando $x = 0$ temos
            \begin{align*}
                0 &= y(0) = C + D \Rightarrow C = - D.
            \end{align*}
            Quando $x = \pi$ temos
            \begin{align*}
                0 &= y(\pi) = C \left[ \exp\left( (3 + k) \pi 2^{-1} \right) - \exp\left( (3 - k) \pi 2^{-1} \right) \right] \\
                &= C \exp(3/2) \left[ \exp(k \pi 2^{-1}) - \exp(- k \pi 2^{-1}) \right] \Rightarrow C = 0.
            \end{align*}
            Portanto, para $-3 - 12 \lambda > 0$ temos apenas a solu\c{c}\~{a}o trivial.


            Para $-3 - 12 \lambda < 0$ temos
            \begin{align*}
                y(x) = E \exp\left( (3 + k i) x 2^{-1} \right) + F \exp\left( (3 - k i) x 2^{-1} \right).
            \end{align*}
            Quando $x = 0$ temos
            \begin{align*}
                0 &= y(0) = E + F \Rightarrow E = - F.
            \end{align*}
            Quando $x = \pi)$ temos
            \begin{align*}
                0 &= y(\pi) = E \exp(3 \pi 2^{-1}) \left[ \exp(k i \pi 2^{-1}) - \exp(- k i \pi 2^{-1}) \right] \\
                &= E \exp(3 \pi 2^{-1}) \left[ \cos(k \pi 2^{-1}) + i \sin(k \pi 2^{-1}) - \cos(k \pi 2^{-1}) + i \sin(k \pi 2^{-1}) \right] \\
                &= \sin(k \pi 2^{-1}) \Rightarrow k = 2n.
            \end{align*}
            Portanto, para $-3 - 12 \lambda < 0$ temos, al\'{e}m da solu\c{c}\~{a}o trivial, $y(x) = \exp(3 x 2^{-1}) \sin(n x)$ onde $k = 2n = \sqrt{- 3 - 12 \lambda}$.

            Assim sendo, temos, al\'{e}m da solu\c{c}\~{a}o trivial, $y(x) = \exp(3 x 2^{-1}) \sin(n x)$, onde $k = 2n = \sqrt{- 3 - 12 \lambda}$, quando $-3 - 12 \lambda < 0$.
        \end{solution}

        \part[P3 de 2006] $\begin{cases}
            \frac{\mathrm{d}}{\mathrm{d}x} \left[ (2 + x)^2 \frac{\mathrm{d}y}{\mathrm{d}x} \right] = -\lambda y, & -1 < x < 1, \\
            y(-1) = 0, \\
            y(1) = 0.
        \end{cases}$
        \begin{solution}
            Temos que
            \begin{align*}
                \frac{\mathrm{d}}{\mathrm{d}x}\left[ (2 + x)^2 \frac{\mathrm{d}y}{\mathrm{d}x} \right] &= (2 + x)^2 \frac{\mathrm{d}^2y}{\mathrm{d}x^2} + 2 (2 + x) \frac{\mathrm{d}y}{\mathrm{d}x} = -\lambda y
            \end{align*}
            \'{e} uma equa\c{c}\~{a}o de Guler. Logo, $y = (2 + x)^r$ e portanto
            \begin{align*}
                r (r - 1) + 2 r + \lambda = 0
            \end{align*}
            que implica em $r = \left( -1 \pm \sqrt{1 - 4 \lambda} \right) 2^{-1}$.

            Para $1 - 4 \lambda > 0$ temos que $r_1 = \left( -1 + \sqrt{1 - 4 \lambda} \right) 2^{-1}$ e $r_2 = \left( -1 - \sqrt{1 - 4 \lambda} \right) 2^{-1}$ e portanto
            \begin{align*}
                y(x) &= A_1 (2 + x)^{r_1} + A_2 (2 + x)^{r_2}.
            \end{align*}
            Quando $x = -1$,
            \begin{align*}
                0 = y(-1) = A_1 + A_2 \Rightarrow A_1 = - A_2.
            \end{align*}
            Quando $x = 1$,
            \begin{align*}
                0 = y(1) = A_1 3^{r_1} + A_2 3^{r_2} = A_1 (3^{r_1} - 3^{r_2}) \Rightarrow A_1 = 0 \Rightarrow = 0.
            \end{align*}
            Assim concluimos que para $1 - 4 \lambda > 0$ temos apenas a solu\c{c}\~{a}o trivial.

            Para $1 - 4 \lambda = 0$ temos que $r_1 = r_2 = -1/2$ e portanto
            \begin{align*}
                y(x) = A_3 (2 + x)^{-1/2} + A_4 (2 + x)^{-1/2} \ln(2 + x).
            \end{align*}
            Quando $x = -1$,
            \begin{align*}
                0 = y(-1) = A_3 \Rightarrow A_3 = 0.
            \end{align*}
            Quando $x = 1$,
            \begin{align*}
                0 = y(1) = A_4 3^{-1/2} \ln(3) \Rightarrow A_4 = 0.
            \end{align*}
            Assim concluimos que para $1 - 4 \lambda = 0$ temos apenas a solu\c{c}\~{a}o trivial.

            Para $1 - 4 \lambda < 0$ fazemos $1 - 4 \lambda = - k^2$, $k > 0$. Temos ent\~{a}o que $r = \left( -1 \pm i k \right) 2^{-1}$ e portanto
            \begin{align*}
                y(x) &= B_1 (x + 2)^{-1/2 + i k 2^{-1}} + B_2 (x + 2)^{-1/2 - i k 2^{-1}} \\
                &= (x + 2)^{-1/2} \left[ B_1 \exp(i k 2^{-1} \ln(x + 2)) + B_2 \exp(-i k 2^{-1} \ln(x + 2)) \right] \\
                &= (x + 2)^{-1/2} \left[ C_1 \cos(k 2^{-1} \ln(x + 2)) + C_2 \sin(k 2^{-1} \ln(x + 2)) \right].
            \end{align*}
            Quando $x = -1$,
            \begin{align*}
                0 = y(-1) = C_1 \Rightarrow C_1 = 0.
            \end{align*}
            Quando $x = 1$,
            \begin{align*}
                0 = y(1) = 3^{-1/2} C_2 \sin(k 2^{-1} \ln(3)) \Rightarrow k 2^{-1} \ln(3) = n \pi, n = 1, 2, \ldots.
            \end{align*}
            Logo,
            \begin{align*}
                1 - 4 \lambda = -\left( 2 n \pi / \ln(3) \right)^2
            \end{align*}
            que implica em
            \begin{align*}
                \lambda_n = 1/4 + \left( n \pi / \ln(3) \right)^2, n = 1, 2, \ldots.
            \end{align*}
            Assim concluimos que para $1 - 4 \lambda < 0$ temos como solu\c{c}\~{a}o
            \begin{align*}
                y_n(x) = (2 + x)^{-1/2} \sin\left( n \pi (\ln(3))^{-1} \ln(2 + x) \right).
            \end{align*}
        \end{solution}

        \part $\begin{cases}
            x^2 y'' + x y' + \lambda y = 0, & 0 < x < 1, \\
            y(1) = 0, \\
            \lim_{x \to 0^+} | y(x) | < \infty.
        \end{cases}$
        \begin{solution}
            Para $y(x) = x^r$ temos que
            \begin{align*}
                r^2 + \lambda = 0 \Rightarrow r = \pm i \sqrt{\lambda}.
            \end{align*}
            Para $\lambda = 0$, temos que
            \begin{align*}
                y = A x^0 + B x^0 \ln(x) = A + B \ln(x).
            \end{align*}
            Quando $x = 1$ temos
            \begin{align*}
                0 &= y(1) = A + B \ln(1) \Rightarrow A = 0.
            \end{align*}
            Quando $x \to 0^+$ temos
            \begin{align*}
                \lim_{x \to 0^+} | y(x) | < \infty \Rightarrow \lim_{x \to 0^+} | B \ln(x) | < \infty.
            \end{align*}
            Portanto, para $\lambda = 0$ n\~{a}o existe solu\c{c}\~{a}o.

            Para $\lambda > 0$, fazendo $\lambda = k^2$ com $k > 0$, temos que
            \begin{align*}
                y(x) &= \bar{A} x^{i k} + \bar{B} x^{-i k} \\
                &= \bar{A} \exp(i k \ln(x)) + \bar{B} \exp(-i k \ln(x)) \\
                &= A \cos(k \ln(x)) + B \sin(k \ln(x)).
            \end{align*}
            Quando $x = 1$ temos
            \begin{align*}
                0 = y(1) = A \cos(k \ln(x)) + B \sin(k \ln(x)) \Rightarrow A = 0.
            \end{align*}
            Quando $x \to 0^+$ temos
            \begin{align*}
                \lim_{x \to 0^+} | y(x) | < \infty \Rightarrow \lim_{x \to 0^+} | B \sin(k \ln(x)) | < \infty.
            \end{align*}
            Portanto, para $\lambda > 0$ temos que $-1 \leq | B \sin(k \ln(x))  \leq 1$ que implica em $y(x) = \sin(k \ln(x))$, $k \geq 1$.

            Para $\lambda = - k^2$, fazendo $\lambda = - k^2$ com $k > 0$, temos que
            \begin{align*}
                y(x) = A x^k + B x^{-k}.
            \end{align*}
            Quando $x = 1$ temos
            \begin{align*}
                0 = y(1) = A + B \Rightarrow A = - B.
            \end{align*}
            Quando $x \to 0^+$ temos
            \begin{align*}
                \lim_{x \to 0^+} | y(x) | = \lim_{x \to 0^+} | x^k - x^{-k} | = -\infty.
            \end{align*}
            Portanto, para $\lambda < 0$ n\~{a}o existe solu\c{c}\~{a}o.

            Assim sendo, temos como solu\c{c}\~{a}o apenas $y(x) = \sin(k \ln(x))$ com $k \geq 1$ quando $\lambda > 0$.
        \end{solution}

        \part $\begin{cases}
            y'' + \lambda y = 0, & 0 < x < \infty, \\
            y(0) = 0, \\
            \lim_{x \to \infty} | y(x) | < \infty.
        \end{cases}$
        \begin{solution}
            Para $\lambda = 0$, temos que
            \begin{align*}
                y'' = 0 \Rightarrow y = A + B x.
            \end{align*}
            Quando $x = 0$ temos
            \begin{align*}
                0 = y(0) = A.
            \end{align*}
            Quando $x \to \infty$ temos
            \begin{align*}
                \lim_{x \to \infty} | y(x) | = \lim_{x \to \infty} | B x | = \infty.
            \end{align*}
            Portanto, para $\lambda = 0$ n\~{a}o existe solu\c{c}\~{a}o.

            Para $\lambda > 0$, fazendo $\lambda = k^2$ com $k > 0$, temos que
            \begin{align*}
                y'' + k^2 y = 0 \Rightarrow y = C \cos(k x) + D \sin(k x).
            \end{align*}
            Quando $x = 0$ temos
            \begin{align*}
                0 = y(0) = C.
            \end{align*}
            Quando $x \to \infty$ temos
            \begin{align*}
                \lim_{x \to \infty} | y(x) | = \lim_{x \to \infty} | D \sin(k x) | < \infty.
            \end{align*}
            Portanto, para $\lambda > 0$ temos $y(x) = \sin(k x)$ com $k > 0$.

            Para $\lambda < 0$ temos
            \begin{align*}
                y'' - k^2 y = 0 \Rightarrow y = E \cosh(k x) + F \sinh(k x).
            \end{align*}
            Quando $x = 0$ temos
            \begin{align*}
                0 = y(0) = E.
            \end{align*}
            Quando $x \to \infty$ temos
            \begin{align*}
                \lim_{x \to \infty} | y(x) | = \lim_{x \to \infty} | F \sinh(k x) | = \infty.
            \end{align*}
            Portanto, para $\lambda < 0$ n\~{a}o existe solu\c{c}\~{a}o.

            Assim sendo, temos como solu\c{c}\~{a}o apenas $y(x) = \sin(k x)$ com $k > 0$ quando $\lambda > 0$.
        \end{solution}

        \part[P2 de 2012, E de 2012] $\begin{cases}
            x y'' + y' + \lambda x y = 0, & 0 < x < 1, \\
            y(1) = 0, \\
            \lim_{x \to 0^+} | y(x) | < \infty.
        \end{cases}$
        \begin{solution}
            Multiplicando a equa\c{c}\~{a}o diferencial por $x$ temos
            \begin{align*}
                x^2 y'' + x y' + \lambda x^2 y &= 0
            \end{align*}
            que \'{e} uma equa\c{c}\~{a}o de Bessel de ordem zero.

            Para $\lambda = 0$ temos
            \begin{align*}
                x y'' + y' = (x y')' = 0 \Rightarrow y(x) = c_1 \ln(x) + c_2.
            \end{align*}
            Quando $x \to 0^+$ temos
            \begin{align*}
                \lim_{x \to 0^+} | y(x) | < \infty \to x_1 = 0.
            \end{align*}
            Quando $x = 1$ temos
            \begin{align*}
                0 &= y(1) = c_2.
            \end{align*}
            Portanto, quando $\lambda = 0$ existe apenas a solu\c{c}\~{a}o trivial.

            Para $\lambda > 0$, fazendo $\lambda = k^2$ com $k > 0$, temos
            \begin{align*}
                x^2 y'' + x y' + k^2 x^2 y = 0 \Rightarrow y = A J_0(k x) + B Y_0(k x)
            \end{align*}
            onde $J_0(k x)$ \'{e} a equa\c{c}\~{a}o de Bessel de primeira esp\'{e}cie e ordem zero e $Y_0(k x)$ \'{e} a equa\c{c}\~{a}o de Bessel de segunda esp\'{e}cie e ordem zero. Quando $x \to 0^+$ temos que
            \begin{align*}
                \lim_{x \to 0^+} | y(x) | < \infty \Rightarrow B = 0
            \end{align*}
            pois $\lim_{x \to 0^+} y(x) \to -\infty$. Quando $x = 1$ temos que
            \begin{align*}
                0 = y(1) = A J_0(k) \Rightarrow J_0(k) = 0 \Rightarrow k = \alpha_{0n}
            \end{align*}
            onde $\alpha_{0n}$ \'{e} o $n$-\'{e}simo zero de $J_0(x)$ para $n = 1, 2, \ldots$. Portanto, quando $\lambda > 0$ temos como solu\c{c}\~{a}o o autopar $(\alpha_{0n}^2, J_0(\alpha_{0n} x)$.

            Para $\lambda < 0$, fazendo $\lambda = - k^2$ com $k > 0$, temos
            \begin{align*}
                x^2 y'' + x y' - k^2 x^2 y = 0 \Rightarrow y = A I_0(k x) + B K_0(k x),
            \end{align*}
            onde $I_0(k x)$ \'{e} a fun\c{c}\~{a}o de Bessel modificada de primeira esp\'{e}cie e ordem zero e $K_0(k x)$ \'{e} a fun\c{c}\~{a}o de Bessel modificada de segunda esp\'{e}cie e ordem zero. Quando $x \to 0^+$ temos
            \begin{align*}
                \lim_{x \to 0^+} | y(x) | < \infty \Rightarrow B = 0
            \end{align*}
            pois $\lim_{x \to 0^+} K_0(x) \to \infty$. Quando $x = 1$ temos
            \begin{align*}
                0 = y(1) = A I_0(k) \Rightarrow A = 0
            \end{align*}
            pois $I_0(k) \neq 0$. Portanto, quando $\lambda < 0$ existe apenas a solu\c{c}\~{a}o trivial.

            Assim sendo, al\'{e}m da solu\c{c}\~{a}o trivial temos o autopar $(\alpha_{0n}^2, J_0(\alpha_{0n} x)$.
        \end{solution}
        
        \part $\begin{cases}
            \frac{\mathrm{d}}{\mathrm{d}x} \left[ (x - a) (b - x) \frac{\mathrm{d}y}{\mathrm{d}x} \right] + \lambda y = 0, & a < x < b, \\
            \lim_{x \to a} | y(x) | < \infty, \\
            \lim_{x \to b} | y(x) | < \infty.
        \end{cases}$
        \begin{solution}
            Fazendo $x = a + t \left( b - a \right)$ temos que
            \begin{align*}
                t &= \left( x - a \right) / \left( b - a \right), \\
                x &= a \to t = 0, \\
                x &= b \to t = 1, \\
                \frac{\mathrm{d}}{\mathrm{d}x} &= \frac{\mathrm{d}}{\mathrm{d}t} \frac{\mathrm{d}t}{\mathrm{d}x} = \frac{1}{b - a} \frac{\mathrm{d}}{\mathrm{d}t}.
            \end{align*}
            Portanto,
            \begin{align*}
                0 &= \frac{\mathrm{d}^2y}{\mathrm{d}x^2} \left( - x^2 + \left( a + b \right) x - a b \right) + \frac{\mathrm{d}}{\mathrm{d}x} \left( -2 x + (a + b) \right) \mathrm{d}y \\
                \begin{split}
                    &= \frac{\mathrm{d}^2}{\mathrm{d}t^2} \frac{1}{(b - a)^2} \left[ - \left( a + t \left( b - a \right) \right)^2 + \left( a + b \right) \left( a + t \left( b - a \right) \right) - a b \right] \\ &\quad {}+ \frac{\mathrm{d}}{\mathrm{d}t} y \frac{1}{b - a} \left[ -2 a + 2 t \left( b - a \right) + \left( a + b) \right) \right] + \lambda y 
                \end{split} \\
                \begin{split}
                    &= \frac{\mathrm{d}^2 y}{\mathrm{d}t^2} \frac{1}{(b - a)^2} \left[ -a t \left( b - a \right) + b t \left( b - a \right) - t^2 \left( b - a \right)^2 \right] \\ &\quad {}+ \frac{\mathrm{d}y}{\mathrm{d}t} \frac{1}{(b - a)} \left[ -2 t \left( b - a \right) + b - a \right] + \lambda y
                \end{split} \\
                &= t \left( 1 - t \right) \frac{\mathrm{d}^2y}{\mathrm{d}t^2} + \left( 1 - 2 t \right) \frac{\mathrm{d}y}{\mathrm{d}t} + \lambda y.
            \end{align*}
            Para $t = \left( 1 - \bar{x} \right) / 2$ temos
            \begin{align*}
                \left( 1 - \bar{x}^2 \right) y'' - 2 \bar{x} y' + \lambda y &= 0.
            \end{align*}
            Para $\lambda = v \left( v + 1 \right)$ temos
            \begin{align*}
                y &= A P_v (\bar{x}) + B Q_v(\bar{x}) \\
                &= A P_v(1 - 2 t) + B Q_v(1 - 2 t) \\
                &= A \,_2F_1(-v, v + 1, 1; t) + B Q_v(1 - 2 t).
            \end{align*}
            Ent\~{a}o, quando $x \to \pm 1$ temos que $\lim_{x \to \pm 1} | y | < \infty$ e portanto $B = 0$.o

            Como 
            \begin{align*}
                _2F_1(-v, v + 1, 1; t) &= \sum_{n = 0}^\infty \frac{(-v)_n (v + 1)_n}{(1)_n} \frac{t^n}{n!}
            \end{align*}
            e deseja-se que seja infinito para $t = 0$ e $t = 1$ temos que se $v \in \mathbb{N}$ obtemos os polin\^{o}mios de Legendre de modo que satisfaz a condi\c{c}\~{a}o desejada. J\'{a} para o caso geral corresponde temos
            \begin{align*}
                \sum_{n = 0}^\infty \frac{(-v)_n (v + 1)_n}{n!} \frac{1}{n!}
            \end{align*}
            e assim
            \begin{align*}
                \frac{a_{n + 1}}{a_n} &= \frac{(-v)_{n + 1} (v + 1)_{n + 1}}{\left( (n + 1)! \right)^2} \to \frac{(n!)^2}{(-v)_n (v + 1)_n} = \frac{(n - v) (n + 1 - v + 1)}{(n + 1)^2} \to 1.
            \end{align*}
            Portanto,
            \begin{align*}
                _2F_1(-v, v + 1, 1; 1) &= \frac{1}{B(v + 1, v)} \int_0^1 t^v \left( 1 - t \right)^{-v - 1} \left( 1 - t \right)^v \,\mathrm{d}t \\
                &= \frac{1}{B(v + 1, - v)} \int_0^1 t^{v} (1 - t)^-1 \,\mathrm{d}t \\
                &= \frac{B(v + 1, 0)}{B(v + 1, - v)}.
            \end{align*}
            Por fim, $y(t) = P_n(1 - 2t)$ para $\lambda_n = n(n + 1)$ e $y_n(x) = P_n(1 - 2 (x - a)/(b - a))$.
            % TODO Colocar a solu\c{c}\~{a}o acima no formato abaixo.

            % Para $\lambda = 0$, temos que
            % \begin{align*}
            %     y'' = 0 \Rightarrow y = A + B x.
            % \end{align*}
            % Quando $x = 0$ temos
            % \begin{align*}
            %     0 = y(0) = A.
            % \end{align*}
            % Quando $x \to \infty$ temos
            % \begin{align*}
            %     \lim_{x \to \infty} | y(x) | = \lim_{x \to \infty} | B x | = \infty.
            % \end{align*}
            % Portanto, para $\lambda = 0$ n\~{a}o existe solu\c{c}\~{a}o.

            % Para $\lambda > 0$, fazendo $\lambda = k^2$ com $k > 0$, temos que
            % \begin{align*}
            %     y'' + k^2 y = 0 \Rightarrow y = C \cos(k x) + D \sin(k x).
            % \end{align*}
            % Quando $x = 0$ temos
            % \begin{align*}
            %     0 = y(0) = C.
            % \end{align*}
            % Quando $x \to \infty$ temos
            % \begin{align*}
            %     \lim_{x \to \infty} | y(x) | = \lim_{x \to \infty} | D \sin(k x) | < \infty.
            % \end{align*}
            % Portanto, para $\lambda > 0$ temos $y(x) = \sin(k x)$ com $k > 0$.

            % Para $\lambda < 0$ temos
            % \begin{align*}
            %     y'' - k^2 y = 0 \Rightarrow y = E \cosh(k x) + F \sinh(k x).
            % \end{align*}
            % Quando $x = 0$ temos
            % \begin{align*}
            %     0 = y(0) = E.
            % \end{align*}
            % Quando $x \to \infty$ temos
            % \begin{align*}
            %     \lim_{x \to \infty} | y(x) | = \lim_{x \to \infty} | F \sinh(k x) | = \infty.
            % \end{align*}
            % Portanto, para $\lambda < 0$ n\~{a}o existe solu\c{c}\~{a}o.

            % Assim sendo, temos como solu\c{c}\~{a}o apenas $y(x) = \sin(k x)$ com $k > 0$ quando $\lambda > 0$.
        \end{solution}

        \part $\begin{cases}
            y'' - y = 0, & 0 < x < 1, \\
            y(0) = 0, \\
            y(1) = 1.
        \end{cases}$
        \begin{solution}
            Temos que a solu\c{c}\~{a}o \'{e} da forma
            \begin{align*}
                y(x) &= A \exp(x) + B \exp(-x).
            \end{align*}
            Quando $x = 0$ temos
            \begin{align*}
                0 &= y(0) = A + B \to B = - A.
            \end{align*}
            Quando $x \to 1$ temos
            \begin{align*}
                1 &= y(1) = A \exp(1) + B \exp(-1) = 2 A \sinh(1) \to A = \frac{1}{2 \sinh(1)}.
            \end{align*}
            Portanto, a solu\c{c}\~{a}o \'{e} $y = \sinh(x) / \sinh(1)$.
            % TODO Colocar a solu\c{c}\~{a}o acima no formato abaixo.

            % Para $\lambda = 0$, temos que
            % \begin{align*}
            %     y'' = 0 \Rightarrow y = A + B x.
            % \end{align*}
            % Quando $x = 0$ temos
            % \begin{align*}
            %     0 = y(0) = A.
            % \end{align*}
            % Quando $x \to \infty$ temos
            % \begin{align*}
            %     \lim_{x \to \infty} | y(x) | = \lim_{x \to \infty} | B x | = \infty.
            % \end{align*}
            % Portanto, para $\lambda = 0$ n\~{a}o existe solu\c{c}\~{a}o.

            % Para $\lambda > 0$, fazendo $\lambda = k^2$ com $k > 0$, temos que
            % \begin{align*}
            %     y'' + k^2 y = 0 \Rightarrow y = C \cos(k x) + D \sin(k x).
            % \end{align*}
            % Quando $x = 0$ temos
            % \begin{align*}
            %     0 = y(0) = C.
            % \end{align*}
            % Quando $x \to \infty$ temos
            % \begin{align*}
            %     \lim_{x \to \infty} | y(x) | = \lim_{x \to \infty} | D \sin(k x) | < \infty.
            % \end{align*}
            % Portanto, para $\lambda > 0$ temos $y(x) = \sin(k x)$ com $k > 0$.

            % Para $\lambda < 0$ temos
            % \begin{align*}
            %     y'' - k^2 y = 0 \Rightarrow y = E \cosh(k x) + F \sinh(k x).
            % \end{align*}
            % Quando $x = 0$ temos
            % \begin{align*}
            %     0 = y(0) = E.
            % \end{align*}
            % Quando $x \to \infty$ temos
            % \begin{align*}
            %     \lim_{x \to \infty} | y(x) | = \lim_{x \to \infty} | F \sinh(k x) | = \infty.
            % \end{align*}
            % Portanto, para $\lambda < 0$ n\~{a}o existe solu\c{c}\~{a}o.

            % Assim sendo, temos como solu\c{c}\~{a}o apenas $y(x) = \sin(k x)$ com $k > 0$ quando $\lambda > 0$.
        \end{solution}

        \part $\begin{cases}
            y'' + 4 y' + 7y = 0, & 0 < x < 1, \\
            y(0) = 0, \\
            y'(1) = 1.
        \end{cases}$
        \begin{solution}
            Temos que a solu\c{c}\~{a}o \'{e} da forma $y(x) = \exp(r x)$. Portanto,
            \begin{align*}
                r^2 + 4 r + 7 = 0 \to r = \left( -4 \pm \sqrt{16 - 28} \right) / 2 = -2 \pm i \sqrt{3}
            \end{align*}
            e consequentemente
            \begin{align*}
                y(x) &= \exp(-2 x) \left( A \cos(\sqrt{3} x) + B \sin(\sqrt{3} x) \right).
            \end{align*}
            Quando $x = 0$ temos
            \begin{align*}
                0 &= y(0) = A \to y = \exp(-2 x) B \sin(\sqrt{3} x).
            \end{align*}
            Quando $x = 1$ temos que
            \begin{align*}
                y'(x) &= -2 \exp(-2 x) B \sin(\sqrt{3} x) + \sqrt{3} \exp(- 2 x) B \cos(\sqrt{3} x), \\
                1 &= y'(1) = B \left[ -2 \exp(-2) \sin(\sqrt{3}) + \sqrt{3}\exp(-2) \cos(\sqrt{3}) \right] \\
                &= \to B = \exp(2) / \left( -2 \sin(\sqrt{3} + \sqrt{3} \cos(\sqrt{3}) \right).
            \end{align*}
            Portanto, a solu\c{c}\~{a}o \'{e} $y = \exp(2) \left( -2 \sin(\sqrt{3}) + \sqrt{3} \cos(\sqrt{3}) \right)^{-1} \exp(-2x) \sin(\sqrt{3} x)$.
            % TODO Colocar a solu\c{c}\~{a}o acima no formato abaixo.

            % Para $\lambda = 0$, temos que
            % \begin{align*}
            %     y'' = 0 \Rightarrow y = A + B x.
            % \end{align*}
            % Quando $x = 0$ temos
            % \begin{align*}
            %     0 = y(0) = A.
            % \end{align*}
            % Quando $x \to \infty$ temos
            % \begin{align*}
            %     \lim_{x \to \infty} | y(x) | = \lim_{x \to \infty} | B x | = \infty.
            % \end{align*}
            % Portanto, para $\lambda = 0$ n\~{a}o existe solu\c{c}\~{a}o.

            % Para $\lambda > 0$, fazendo $\lambda = k^2$ com $k > 0$, temos que
            % \begin{align*}
            %     y'' + k^2 y = 0 \Rightarrow y = C \cos(k x) + D \sin(k x).
            % \end{align*}
            % Quando $x = 0$ temos
            % \begin{align*}
            %     0 = y(0) = C.
            % \end{align*}
            % Quando $x \to \infty$ temos
            % \begin{align*}
            %     \lim_{x \to \infty} | y(x) | = \lim_{x \to \infty} | D \sin(k x) | < \infty.
            % \end{align*}
            % Portanto, para $\lambda > 0$ temos $y(x) = \sin(k x)$ com $k > 0$.

            % Para $\lambda < 0$ temos
            % \begin{align*}
            %     y'' - k^2 y = 0 \Rightarrow y = E \cosh(k x) + F \sinh(k x).
            % \end{align*}
            % Quando $x = 0$ temos
            % \begin{align*}
            %     0 = y(0) = E.
            % \end{align*}
            % Quando $x \to \infty$ temos
            % \begin{align*}
            %     \lim_{x \to \infty} | y(x) | = \lim_{x \to \infty} | F \sinh(k x) | = \infty.
            % \end{align*}
            % Portanto, para $\lambda < 0$ n\~{a}o existe solu\c{c}\~{a}o.

            % Assim sendo, temos como solu\c{c}\~{a}o apenas $y(x) = \sin(k x)$ com $k > 0$ quando $\lambda > 0$.
        \end{solution}

        \part $\begin{cases}
            y'' - 6 y' + 25 y = 0, & 0 < x < \pi/4, \\
            y'(0) = 1, \\
            y(\pi/4) = 0.
        \end{cases}$
        \begin{solution}
            Temos que a solu\c{c}\~{a}o \'{e} da forma $y(x) = \exp(r x)$. Portanto,
            \begin{align*}
                r^2 - 6 r + 25 = 0 \to r = \left( 6 \pm \sqrt{36 - 100} \right) / 2 = \left( 6 \pm 8 i \right) / 2 = 3 \pm 4 i
            \end{align*}
            e consequentemente
            \begin{align*}
                y(x) = \exp(3 x) \left( A \cos(4x) + B \sin(4x) \right).
            \end{align*}
            Quando $x = \pi/4$ temos
            \begin{align*}
                0 &= y(\pi/4) = \exp(3 \pi / 4) \left( A \cos(\pi) + B \sin(\pi) \right) \to A = 0.
            \end{align*}
            Quando $x = 0$ temos
            \begin{align*}
                y'(x) &= B \left( \exp(3 x) 3 \sin(4 x) + \exp(3 x) 4 \cos(4 x) \right), \\
                1 &= y'(0) = B(4) \to B = 1/4.
            \end{align*}
            Portanto, a solu\c{c}\~{a}o \'{e} $y(x) = \left( 1/4 \right) \exp(3 x) \sin(4x)$. 
            % TODO Colocar a solu\c{c}\~{a}o acima no formato abaixo.

            % Para $\lambda = 0$, temos que
            % \begin{align*}
            %     y'' = 0 \Rightarrow y = A + B x.
            % \end{align*}
            % Quando $x = 0$ temos
            % \begin{align*}
            %     0 = y(0) = A.
            % \end{align*}
            % Quando $x \to \infty$ temos
            % \begin{align*}
            %     \lim_{x \to \infty} | y(x) | = \lim_{x \to \infty} | B x | = \infty.
            % \end{align*}
            % Portanto, para $\lambda = 0$ n\~{a}o existe solu\c{c}\~{a}o.

            % Para $\lambda > 0$, fazendo $\lambda = k^2$ com $k > 0$, temos que
            % \begin{align*}
            %     y'' + k^2 y = 0 \Rightarrow y = C \cos(k x) + D \sin(k x).
            % \end{align*}
            % Quando $x = 0$ temos
            % \begin{align*}
            %     0 = y(0) = C.
            % \end{align*}
            % Quando $x \to \infty$ temos
            % \begin{align*}
            %     \lim_{x \to \infty} | y(x) | = \lim_{x \to \infty} | D \sin(k x) | < \infty.
            % \end{align*}
            % Portanto, para $\lambda > 0$ temos $y(x) = \sin(k x)$ com $k > 0$.

            % Para $\lambda < 0$ temos
            % \begin{align*}
            %     y'' - k^2 y = 0 \Rightarrow y = E \cosh(k x) + F \sinh(k x).
            % \end{align*}
            % Quando $x = 0$ temos
            % \begin{align*}
            %     0 = y(0) = E.
            % \end{align*}
            % Quando $x \to \infty$ temos
            % \begin{align*}
            %     \lim_{x \to \infty} | y(x) | = \lim_{x \to \infty} | F \sinh(k x) | = \infty.
            % \end{align*}
            % Portanto, para $\lambda < 0$ n\~{a}o existe solu\c{c}\~{a}o.

            % Assim sendo, temos como solu\c{c}\~{a}o apenas $y(x) = \sin(k x)$ com $k > 0$ quando $\lambda > 0$.
        \end{solution}

        \part $\begin{cases}
            y'' + 2 y' + y = x, & 0 < x < 2, \\
            y(0) = 0, \\
            y(2) = 3.
        \end{cases}$
        \begin{solution}
            Temos que a solu\c{c}\~{a}o \'{e} da forma $y(x) = \exp(r x)$. Portanto,
            \begin{align*}
                r^2 - 6 r + 25 = 0 \to r = \left( 6 \pm \sqrt{36 - 100} \right) / 2 = \left( 6 \pm 8 i \right) / 2 = 3 \pm 4 i
            \end{align*}
            e consequentemente
            \begin{align*}
                y(x) &= \exp(3 x) \left( A \cos(4 x) + B \sin(4 x) \right).
            \end{align*}
            Uma solu\c{c}\~{a}o particular \'{e} da forma $\bar{A} x + \bar{B}$ e portanto
            \begin{align*}
                2 (\bar{A}) + \bar{A} + \bar{B} = x \to \bar{A} = 1, \bar{B} = -2.
            \end{align*}
            Logo, a solu\c{c}\~{a}o \'{e} da forma
            \begin{align*}
                y(x) &= x - 2 + \exp(3 x) \left( A \cos(4 x) + B \sin(4 x) \right).
            \end{align*}
            Quando $x = 0$ temos
            \begin{align*}
                0 &= y(0) = -2 + A \to A = 2.
            \end{align*}
            Quando $x = 2$ temos
            \begin{align*}
                3 &= y(2) = \exp(6) \left( 2 \cos(8) + B \sin(8) \right) \to B = - 2 \cot(8).
            \end{align*}
            Portanto, a solu\c{c}\~{a}o \'{e} $y = x - 2 + \exp(3 x) \left( 2 \cos(4 x) - 2 \cot(8) \sin(4 x) \right)$.
            % TODO Colocar a solu\c{c}\~{a}o acima no formato abaixo.

            % Para $\lambda = 0$, temos que
            % \begin{align*}
            %     y'' = 0 \Rightarrow y = A + B x.
            % \end{align*}
            % Quando $x = 0$ temos
            % \begin{align*}
            %     0 = y(0) = A.
            % \end{align*}
            % Quando $x \to \infty$ temos
            % \begin{align*}
            %     \lim_{x \to \infty} | y(x) | = \lim_{x \to \infty} | B x | = \infty.
            % \end{align*}
            % Portanto, para $\lambda = 0$ n\~{a}o existe solu\c{c}\~{a}o.

            % Para $\lambda > 0$, fazendo $\lambda = k^2$ com $k > 0$, temos que
            % \begin{align*}
            %     y'' + k^2 y = 0 \Rightarrow y = C \cos(k x) + D \sin(k x).
            % \end{align*}
            % Quando $x = 0$ temos
            % \begin{align*}
            %     0 = y(0) = C.
            % \end{align*}
            % Quando $x \to \infty$ temos
            % \begin{align*}
            %     \lim_{x \to \infty} | y(x) | = \lim_{x \to \infty} | D \sin(k x) | < \infty.
            % \end{align*}
            % Portanto, para $\lambda > 0$ temos $y(x) = \sin(k x)$ com $k > 0$.

            % Para $\lambda < 0$ temos
            % \begin{align*}
            %     y'' - k^2 y = 0 \Rightarrow y = E \cosh(k x) + F \sinh(k x).
            % \end{align*}
            % Quando $x = 0$ temos
            % \begin{align*}
            %     0 = y(0) = E.
            % \end{align*}
            % Quando $x \to \infty$ temos
            % \begin{align*}
            %     \lim_{x \to \infty} | y(x) | = \lim_{x \to \infty} | F \sinh(k x) | = \infty.
            % \end{align*}
            % Portanto, para $\lambda < 0$ n\~{a}o existe solu\c{c}\~{a}o.

            % Assim sendo, temos como solu\c{c}\~{a}o apenas $y(x) = \sin(k x)$ com $k > 0$ quando $\lambda > 0$.
        \end{solution}

        \part $\begin{cases}
            \left( x^2 y' \right)' + \lambda x^{-2}y = 0, & 1 < x < 2, \\
            y(1) = 0,\\
            y(2) = 0.
        \end{cases}$
        \begin{solution}
            A equa\c{c}\~{a}o diferencial corresponde a
            \begin{align*}
                x^2 y'' + 2 x y' + \lambda / x^2 y &= 0, \\
                y'' + 2 / x y' + \lambda / x^4 y &= 0
            \end{align*}
            que possue ponto singular irregular em $x = 0$.

            Fazendo $t = 1/x$ temos
            \begin{align*}
                \frac{\mathrm{d}}{\mathrm{d}x} &= \frac{\mathrm{d}t}{\mathrm{d}x} \frac{\mathrm{d}}{\mathrm{d}t} = -t^2 \frac{\mathrm{d}}{\mathrm{d}t}, \\
                \frac{\mathrm{d}^2}{\mathrm{d}x^2} &= t^4 \frac{\mathrm{d}^2}{\mathrm{d}t^2} + 2 t^3 \frac{\mathrm{d}}{\mathrm{d}t}.
            \end{align*}
            Substituindo temos
            \begin{align*}
                t^4 \frac{\mathrm{d}^2y}{\mathrm{d}t^2} + 2 t^3 \frac{\mathrm{d}y}{\mathrm{d}x} + 2 t \left( - t^2 \frac{\mathrm{d}y}{\mathrm{d}t} \right) + \lambda t^4 y = 0
            \end{align*}
            e portanto re-escrevemos o problema de Sturm-Liouville como
            \begin{align*}
                \begin{cases}
                    \frac{\mathrm{d}^2y}{\mathrm{d}t^2} + \lambda y = 0, 1/2 < t < 1, \\
                    y(1/2) = 0, \\
                    y(1) = 0.
                \end{cases}
            \end{align*}

            Para $\lambda > 0$, fazendo $\lambda = k^2$ com $k > 0$, temos que
            \begin{align*}
                y = A \cos(k t) + B \sin(k t).
            \end{align*}
            Quando $x = 1$ temos
            \begin{align*}
                0 &= y(1) = A \cos(k) + B \sin(k) \to A = - B \sin(k) / \cos(k).
            \end{align*}
            Quando $x = 1/2$ temos
            \begin{align*}
                0 &= y(1/2) = A \cos(k/2) + B \sin(k/2) \\
                &= B \sin(k) \cos(k/2) / \cos(k) + B \sin(k/2) \\
                &= \sin(k) \cos(k/2) + \cos(k) \sin(k/2) && B \neq 0 \\
                &= \sin(k - k/2) = \sin(k/2).
            \end{align*}
            Logo, $k/2 = n \pi$, $n = 1, 2, \ldots$ e
            \begin{align*}
                y = - B \sin(k) \cos(k t) / \cos(k) + B \sin(k t) = \bar{B} \sin(k (1 - t)).
            \end{align*}
            Desfazendo a mudan\c{c}a de vari\'{a}vel temos
            \begin{align*}
                y_n(x) &= \sin(2 n \pi / x), n = 1, 2, \ldots, \\
                \lambda_n &= 4 n^2 \pi^2.
            \end{align*}
            % TODO Colocar a solu\c{c}\~{a}o acima no formato abaixo.

            % Para $\lambda = 0$, temos que
            % \begin{align*}
            %     y'' = 0 \Rightarrow y = A + B x.
            % \end{align*}
            % Quando $x = 0$ temos
            % \begin{align*}
            %     0 = y(0) = A.
            % \end{align*}
            % Quando $x \to \infty$ temos
            % \begin{align*}
            %     \lim_{x \to \infty} | y(x) | = \lim_{x \to \infty} | B x | = \infty.
            % \end{align*}
            % Portanto, para $\lambda = 0$ n\~{a}o existe solu\c{c}\~{a}o.

            % Para $\lambda > 0$, fazendo $\lambda = k^2$ com $k > 0$, temos que
            % \begin{align*}
            %     y'' + k^2 y = 0 \Rightarrow y = C \cos(k x) + D \sin(k x).
            % \end{align*}
            % Quando $x = 0$ temos
            % \begin{align*}
            %     0 = y(0) = C.
            % \end{align*}
            % Quando $x \to \infty$ temos
            % \begin{align*}
            %     \lim_{x \to \infty} | y(x) | = \lim_{x \to \infty} | D \sin(k x) | < \infty.
            % \end{align*}
            % Portanto, para $\lambda > 0$ temos $y(x) = \sin(k x)$ com $k > 0$.

            % Para $\lambda < 0$ temos
            % \begin{align*}
            %     y'' - k^2 y = 0 \Rightarrow y = E \cosh(k x) + F \sinh(k x).
            % \end{align*}
            % Quando $x = 0$ temos
            % \begin{align*}
            %     0 = y(0) = E.
            % \end{align*}
            % Quando $x \to \infty$ temos
            % \begin{align*}
            %     \lim_{x \to \infty} | y(x) | = \lim_{x \to \infty} | F \sinh(k x) | = \infty.
            % \end{align*}
            % Portanto, para $\lambda < 0$ n\~{a}o existe solu\c{c}\~{a}o.

            % Assim sendo, temos como solu\c{c}\~{a}o apenas $y(x) = \sin(k x)$ com $k > 0$ quando $\lambda > 0$.
        \end{solution}

        \part $\begin{cases}
            x^2 y'' + xy' + \left( \lambda x^2 - 1/4 \right)y = 0, & \pi/2 < x < 3 \pi 2^{-1}, \\
            y(\pi/2) = 0, \\
            y(3 \pi 2^{-1}) = 0.
        \end{cases}$
        \begin{solution}
            Fazendo a seguinte mudan\c{c}a de vari\'{a}vel $x = \lambda^{-1/2} z$ obtemos uma equa\c{c}\~{a}o de Bessel.
            % TODO Terminar de escrever a solu\c{c}\~{a}o e coloc\'{a}-la no formato abaixo.

            % Para $\lambda = 0$, temos que
            % \begin{align*}
            %     y'' = 0 \Rightarrow y = A + B x.
            % \end{align*}
            % Quando $x = 0$ temos
            % \begin{align*}
            %     0 = y(0) = A.
            % \end{align*}
            % Quando $x \to \infty$ temos
            % \begin{align*}
            %     \lim_{x \to \infty} | y(x) | = \lim_{x \to \infty} | B x | = \infty.
            % \end{align*}
            % Portanto, para $\lambda = 0$ n\~{a}o existe solu\c{c}\~{a}o.

            % Para $\lambda > 0$, fazendo $\lambda = k^2$ com $k > 0$, temos que
            % \begin{align*}
            %     y'' + k^2 y = 0 \Rightarrow y = C \cos(k x) + D \sin(k x).
            % \end{align*}
            % Quando $x = 0$ temos
            % \begin{align*}
            %     0 = y(0) = C.
            % \end{align*}
            % Quando $x \to \infty$ temos
            % \begin{align*}
            %     \lim_{x \to \infty} | y(x) | = \lim_{x \to \infty} | D \sin(k x) | < \infty.
            % \end{align*}
            % Portanto, para $\lambda > 0$ temos $y(x) = \sin(k x)$ com $k > 0$.

            % Para $\lambda < 0$ temos
            % \begin{align*}
            %     y'' - k^2 y = 0 \Rightarrow y = E \cosh(k x) + F \sinh(k x).
            % \end{align*}
            % Quando $x = 0$ temos
            % \begin{align*}
            %     0 = y(0) = E.
            % \end{align*}
            % Quando $x \to \infty$ temos
            % \begin{align*}
            %     \lim_{x \to \infty} | y(x) | = \lim_{x \to \infty} | F \sinh(k x) | = \infty.
            % \end{align*}
            % Portanto, para $\lambda < 0$ n\~{a}o existe solu\c{c}\~{a}o.

            % Assim sendo, temos como solu\c{c}\~{a}o apenas $y(x) = \sin(k x)$ com $k > 0$ quando $\lambda > 0$.
        \end{solution}

        \part $\begin{cases}
            y^{(4)} - \lambda y = 0, & 0 < x < 1, \\
            y(0) = y''(0) = y(1) = y''(1) = 0.
        \end{cases}$
        \begin{solution}
            Para $\lambda = 0$, temos que
            \begin{align*}
                y^{(4)} = 0 \Rightarrow y = A x^3 + B x^3 + C x + D.
            \end{align*}
            Quando $x = 0$ temos
            \begin{align*}
                0 &= y(0) = A 0^3 + B 0^3 + C 0 + D \to D = 0, \\
                0 &= y''(0) = 3 (2) A 0 + 3 (2) B \to B = 0.
            \end{align*}
            Quando $x  = 1$ temos
            \begin{align*}
                0 &= y''(1) = 3 (2) A 1 \to A = 0, \\
                0 &= y(1) = C 1 \to C = 0.
            \end{align*}
            Portanto, para $\lambda = 0$ s\'{o} existe a solu\c{c}\~{a}o trivial.
            % TODO Colocar a solu\c{c}\~{a}o no formato abaixo.
            % Para $\lambda > 0$, fazendo $\lambda = k^2$ com $k > 0$, temos que
            % \begin{align*}
            %     y'' + k^2 y = 0 \Rightarrow y = C \cos(k x) + D \sin(k x).
            % \end{align*}
            % Quando $x = 0$ temos
            % \begin{align*}
            %     0 = y(0) = C.
            % \end{align*}
            % Quando $x \to \infty$ temos
            % \begin{align*}
            %     \lim_{x \to \infty} | y(x) | = \lim_{x \to \infty} | D \sin(k x) | < \infty.
            % \end{align*}
            % Portanto, para $\lambda > 0$ temos $y(x) = \sin(k x)$ com $k > 0$.

            % Para $\lambda < 0$ temos
            % \begin{align*}
            %     y'' - k^2 y = 0 \Rightarrow y = E \cosh(k x) + F \sinh(k x).
            % \end{align*}
            % Quando $x = 0$ temos
            % \begin{align*}
            %     0 = y(0) = E.
            % \end{align*}
            % Quando $x \to \infty$ temos
            % \begin{align*}
            %     \lim_{x \to \infty} | y(x) | = \lim_{x \to \infty} | F \sinh(k x) | = \infty.
            % \end{align*}
            % Portanto, para $\lambda < 0$ n\~{a}o existe solu\c{c}\~{a}o.

            % Assim sendo, temos como solu\c{c}\~{a}o apenas $y(x) = \sin(k x)$ com $k > 0$ quando $\lambda > 0$.
            % 
        \end{solution}

        \part $\begin{cases}
            y^{(4)} - \lambda y = 0, & 0 < x < 1, \\
            y(0) = y'(0) = y''(1) = y'''(1) = 0.
        \end{cases}$
        \begin{solution}
            % TODO Escrever a solu\c{c}\~{a}o no formato abaixo.

            % Para $\lambda = 0$, temos que
            % \begin{align*}
            %     y'' = 0 \Rightarrow y = A + B x.
            % \end{align*}
            % Quando $x = 0$ temos
            % \begin{align*}
            %     0 = y(0) = A.
            % \end{align*}
            % Quando $x \to \infty$ temos
            % \begin{align*}
            %     \lim_{x \to \infty} | y(x) | = \lim_{x \to \infty} | B x | = \infty.
            % \end{align*}
            % Portanto, para $\lambda = 0$ n\~{a}o existe solu\c{c}\~{a}o.

            % Para $\lambda > 0$, fazendo $\lambda = k^2$ com $k > 0$, temos que
            % \begin{align*}
            %     y'' + k^2 y = 0 \Rightarrow y = C \cos(k x) + D \sin(k x).
            % \end{align*}
            % Quando $x = 0$ temos
            % \begin{align*}
            %     0 = y(0) = C.
            % \end{align*}
            % Quando $x \to \infty$ temos
            % \begin{align*}
            %     \lim_{x \to \infty} | y(x) | = \lim_{x \to \infty} | D \sin(k x) | < \infty.
            % \end{align*}
            % Portanto, para $\lambda > 0$ temos $y(x) = \sin(k x)$ com $k > 0$.

            % Para $\lambda < 0$ temos
            % \begin{align*}
            %     y'' - k^2 y = 0 \Rightarrow y = E \cosh(k x) + F \sinh(k x).
            % \end{align*}
            % Quando $x = 0$ temos
            % \begin{align*}
            %     0 = y(0) = E.
            % \end{align*}
            % Quando $x \to \infty$ temos
            % \begin{align*}
            %     \lim_{x \to \infty} | y(x) | = \lim_{x \to \infty} | F \sinh(k x) | = \infty.
            % \end{align*}
            % Portanto, para $\lambda < 0$ n\~{a}o existe solu\c{c}\~{a}o.

            % Assim sendo, temos como solu\c{c}\~{a}o apenas $y(x) = \sin(k x)$ com $k > 0$ quando $\lambda > 0$.
        \end{solution}

        \part $\begin{cases}
            (1 - x^2) y'' - 2 x y' + \lambda y = 0, & 0 < x < 1, \\
            y(0) - 0, \\
            \lim_{x \to 1} | y(x) | < \infty.
        \end{cases}$
        \begin{solution}
            % TODO Escrever a solu\c{c}\~{a}o no formato abaixo.

            % Para $\lambda = 0$, temos que
            % \begin{align*}
            %     y'' = 0 \Rightarrow y = A + B x.
            % \end{align*}
            % Quando $x = 0$ temos
            % \begin{align*}
            %     0 = y(0) = A.
            % \end{align*}
            % Quando $x \to \infty$ temos
            % \begin{align*}
            %     \lim_{x \to \infty} | y(x) | = \lim_{x \to \infty} | B x | = \infty.
            % \end{align*}
            % Portanto, para $\lambda = 0$ n\~{a}o existe solu\c{c}\~{a}o.

            % Para $\lambda > 0$, fazendo $\lambda = k^2$ com $k > 0$, temos que
            % \begin{align*}
            %     y'' + k^2 y = 0 \Rightarrow y = C \cos(k x) + D \sin(k x).
            % \end{align*}
            % Quando $x = 0$ temos
            % \begin{align*}
            %     0 = y(0) = C.
            % \end{align*}
            % Quando $x \to \infty$ temos
            % \begin{align*}
            %     \lim_{x \to \infty} | y(x) | = \lim_{x \to \infty} | D \sin(k x) | < \infty.
            % \end{align*}
            % Portanto, para $\lambda > 0$ temos $y(x) = \sin(k x)$ com $k > 0$.

            % Para $\lambda < 0$ temos
            % \begin{align*}
            %     y'' - k^2 y = 0 \Rightarrow y = E \cosh(k x) + F \sinh(k x).
            % \end{align*}
            % Quando $x = 0$ temos
            % \begin{align*}
            %     0 = y(0) = E.
            % \end{align*}
            % Quando $x \to \infty$ temos
            % \begin{align*}
            %     \lim_{x \to \infty} | y(x) | = \lim_{x \to \infty} | F \sinh(k x) | = \infty.
            % \end{align*}
            % Portanto, para $\lambda < 0$ n\~{a}o existe solu\c{c}\~{a}o.

            % Assim sendo, temos como solu\c{c}\~{a}o apenas $y(x) = \sin(k x)$ com $k > 0$ quando $\lambda > 0$.
        \end{solution}
    \end{parts}

    \question[E de 2006] Encontre os autovalores e autofun\c{c}\~{o}es do problema
    \begin{align*}
        \begin{cases}
            x^2 y'' + 3 x y' - y = - \lambda y, & 1 < x < 2, \\
            y(1) = 0, \\
            y(2) = 0.
        \end{cases}
    \end{align*}
    \begin{solution}
        Temos que
        \begin{align*}
            x^2 y'' + 3 x y' + (\lambda - 1) y &= 0
        \end{align*}
        \'{e} uma equa\c{c}\~{a}o de Euler de modo que $y = x^r$ e portanto
        \begin{align*}
            r (r - 1) + 3 r + (\lambda - 1) = 0.
        \end{align*}
        Logo, $r = -1 \pm \sqrt{2 - \lambda}$.

        Para $2 - \lambda > 0$ temos que $r_1 = - 1 + \sqrt{2 - \lambda}$ e $r_2 = -1 - \sqrt{2 - \lambda}$ e portanto
        \begin{align*}
            y(x) = A_1 x^{r_1} + A_2 x^{r_2}.
        \end{align*}
        Quando $x = 1$,
        \begin{align*}
            0 = y(1) = A_1 + A_2.
        \end{align*}
        Quando $x = 2$,
        \begin{align*}
            0 = y(2) = A_1 2^{r_1} + A_2 2^{r_2}.
        \end{align*}
        Pelas duas equa\c{c}\~{a}os anteriores temos que $A_1 = A_2 = 0$ e assim concluimos que para $2 - \lambda > 0$ temos apenas a solu\c{c}\~{a}o trivial.

        Para $2 - \lambda = 0$ temos que $r_1 = r_2 = -1$ e portanto
        \begin{align*}
            y(x) = B_1 x^{-1} + B_2 x^{-1} \ln(x).
        \end{align*}
        Quando $x = 1$,
        \begin{align*}
            0 = y(1) = B_1.
        \end{align*}
        Quando $x = 2$,
        \begin{align*}
            0 = y(2) = B_1 2^{-1} + B_2 2^{-1} \ln(2).
        \end{align*}
        Pelas duas equa\c{c}\~{o}es anteriores temos que $B_1 = B_2 = 0$ e assim concluimos que para $2 - \lambda = 0$ temos apenas a solu\c{c}\~{a}o tivial.

        Para $2 - \lambda < 0$ fazemos $2 - \lambda = - k^2$, $k > 0$. Temos ent\~{a}o que $r_1 = -1 + i k$ e $r_2 = -1 - i k$ e portanto
        \begin{align*}
            y(x) &= C_1 x^{-1 + i k} + C_2 x^{-1 - i k} \\
            &= x^{-1} \left( C_1 x^{i k} + C_2 x^{-i k} \right) \\
            &= x^{-1} \left( C_1 \exp(i k \ln(x)) + C_2 \exp(-i k \ln(x)) \right) \\
            &= x^{-1} \left( C_1' \cos(k \ln(x)) + C_2' \sin(k \ln(X)) \right).
        \end{align*}
        Quando $x = 1$,
        \begin{align*}
            0 = y(1) = 1 \left( C_1' + 0 \right) \Rightarrow C_1' = 0.
        \end{align*}
        Quando $x = 2$,
        \begin{align*}
            0 = y(2) = 2^{-1} C_2' \sin(k \ln(2)) \Rightarrow k \ln(2) = n \pi, n = 1, 2, \ldots.
        \end{align*}
        Assim, os autovalores s\~{a}o $\lambda_n = 2 + k_n^2 = 2 + (n \pi / \ln(2))^2$, $n = 1, 2, \ldots$, e as autofun\c{c}\~{o}es s\~{a}o
        \begin{align*}
            y_n(x) = x^{-1} \sin \left( n \pi (\ln(2))^{-1} \ln(x) \right), n = 1, 2, \ldots.
        \end{align*}
    \end{solution}

    \question[E de 2008] Resolva o seguinte problema de Sturm-Liouville
    \begin{align*}
        \begin{cases}
            y_n''(x) + y_n'(x) = \lambda_n y_n(x), \\
            y_n(0) = 0, \\
            y_n(1) = 0
        \end{cases}
    \end{align*}
    e determine o produto interno do espa\c{c}o de suas auto-fun\c{c}\~{o}es.
    \begin{solution}
        A forma auto-adjunta \'{e}
        \begin{align*}
            \frac{\mathrm{d}}{\mathrm{d}x}\left[ \exp(x) \frac{\mathrm{d}}{\mathrm{d}x}y_n \right] &= \lambda_n \underbrace{\exp(x)}_{\text{fun\c{c}\~{a}o peso}} y_n(x)
        \end{align*}
        e consequentemente a solu\c{c}\~{a}o \'{e} da forma $y = A \exp(k x)$ e o produto interno
        \begin{align*}
            (y_n, y_m) = \int_0^1 \exp(x) y_n y_m \,\mathrm{d}x.
        \end{align*}

        A equa\c{c}\~{a}o caracter\'{i}stica \'{e}
        \begin{align*}
            k^2 + k - \lambda_n = 0 \to k = \frac{-1 \pm \sqrt{1 + 4 \lambda}}{2}.
        \end{align*}

        Quando $1 + 4 \lambda > 0$ temos que a solu\c{c}\~{a}o geral \'{e} da forma
        \begin{align*}
            y = A \exp\left( \frac{-1 + \sqrt{1 + 4\lambda}}{2} x \right) + B \exp\left( \frac{-1 - \sqrt{1 + 4 \lambda}}{2} x \right).
        \end{align*}
        Pelas condi\c{c}\~{o}es de contorno temos que $A = B = 0$.

        Quando $1 + 4 \lambda = 0$ temos que a solu\c{c}\~{a}o geral \'{e} da forma
        \begin{align*}
            y &= A \exp(- x/2) + B x \exp(-x/2) && \text{varia\c{c}\~{a}o de par\^{a}metr}.
        \end{align*}
        Pelas condi\c{c}\~{o}es de contorno temos que $A = B = 0$.

        Quando $1 + 4 \lambda = - w^2 < 0$ a solu\c{c}\~{a}o \'{e} da forma
        \begin{align*}
            y = \exp(-x/2) \left[ A \cos(\omega x) + B \sin(\omega x) \right].
        \end{align*}
        Pelas condi\c{c}\~{o}es de contorno, $A = 0$ e $\sin(\omega) = 0$ que implica em $\omega = n \pi$.

        Logo, a solu\c{c}\~{a}o do problema de Sturm-Liouville \'{e} $(\lambda_n, y_n)$, com $\lambda_n = -(1/4) \left( 1 + n^2 \pi^2 \right)$ e $y_n = \exp(-x/2) \sin(n \pi x)$.
    \end{solution}

    \question[E de 2008] Resolva o seguinte problema de Sturm-Liouville
    \begin{align*}
        \begin{cases}
            y''(x) - 2 y'(x) = \lambda y(x), \\
            y(0) = 0, \\
            y'(1) = 0,
        \end{cases}
    \end{align*}
    e escreva o produto interno do espa\c{c}o das autofun\c{c}\~{o}es.
    \begin{solution}
        Asolu\c{c}\~{a}o geral da eua\c{c}\~{a}o \'{e}
        \begin{align*}
            y(x) = \exp(x) \left( A \exp(\sqrt{1 + \lambda} x) + B \exp(-\sqrt{1 - \lambda} x) \right).
        \end{align*}

        Para $1 + \lambda > 0$ temos que $y(0) = 0$ implica que $A = - B$ e que $y'(1) = 0$ implica
        \begin{align*}
            & \exp(\sqrt{1 + \lambda}) + \exp(-\sqrt{1 - \lambda}) + \sqrt{1 + \lambda} \left( \exp(\sqrt{1 - \lambda}) - \exp(-\sqrt{1 + \lambda}) \right) = 0, \\
            & (1 + \lambda)^{-1/2} = - \tanh(\sqrt{1 + \lambda})
        \end{align*}
        que n\~{a}o possui solu\c{c}\~{a}o diferente da trivial, $A = 0$.

        Para $1 + \lambda = 0$ a solu\c{c}\~{a}o geral \'{e} $y(x) = A \exp(x) + B \exp(x)$. A condi\c{c}\~{a}o de contorno $y(0) = 0$ implica que $A = 0$. A segunda condi\c{c}\~{a}o implica que a \'{u}nica solu\c{c}\~{a}o \'{e} a trivial $A = B = 0$.

        Para $1 + \lambda = - \omega^2 < 0$ temos que $y(x) = \exp(x) \left( A \cos(\omega x) + B \sin(\omega x) \right)$. A primeira condi\c{c}\~{a}o de contorno implica $A = 0$ e a segunda em $\sin(\omega) + \omega \cos(\omega) = 0$ que por sua vez em $- \omega = \tan(\omega)$ que possui infinitas solu\c{c}\~{o}es. 

        A solu\c{c}\~{a}o do problema de Sturm-Liouville ser\'{a}, portanto, $(\lambda_n, \exp(x) \sin(\omega_n x))$, com $\lambda_n = - (1 + w_n^2)$, sendo $\omega_n$ a $n$-\'{e}sima raiz da equa\c{c}\~{a}o transcendental acima.

        A forma auto-adjunta da equa\c{c}\~{a}o \'{e}
        \begin{align*}
            \frac{\mathrm{d}}{\mathrm{d}x}\left( \exp(-2x) y'(x) \right) = \lambda \exp(-2 x) y(x),
        \end{align*}
        de onde se conclui que o produto interno do espa\c{c}o das autofun\c{c}\~{o}es ser\'{a}
        \begin{align*}
            (y_n, y_m) = \int_0^1 y_n(x) y_m(x) \exp(-2x) \,\mathrm{d}x.
        \end{align*}
    \end{solution}

    \question[T4 de 2010] Determine quais das equa\c{c}\~{o}es abaixo possuem solu\c{c}\~{a}o oscilat\'{o}ria (isto \'{e}, com um n\'{u}mero infinito de zeros) no intervalo $(0, \infty)$. Justifique detalhadamente suas respostas.
    \begin{parts}
        \part $y'' + (\cos x + 2) y = 0$,
        \begin{solution}
            % TODO Escrever solu\c{c}\~{a}o.
        \end{solution}

        \part $y'' - \exp(-x) y = 0$,
        \begin{solution}
           % TODO Escrever solu\c{c}\~{a}o. 
        \end{solution}

        \part $y'' + x^{-1} y = 0$.
        \begin{solution}
            % TODO Escrever solu\c{c}\~{a}o.
        \end{solution}
    \end{parts}

    \question[T4 de 2010] Considere a equa\c{c}\~{a}o diferencial $x^2 u'' + \lambda u = 0$ no intervalo $[1, \exp(1)]$.
    \begin{parts}
        \part Determine os autovalores e autofun\c{c}\~{o}es correspondentes para o problema de Sturm-Liouville associados \`{a} EDO acima sujeita \`{a}s condi\c{c}\~{o}es de controno $u(1) = u(\exp(1)) = 0$.
        \begin{solution}
            % TODO Escrever solu\c{c}\~{a}o.
        \end{solution}

        \part Escreva a condi\c{c}\~{a}o de ortogonalidade entre essas autofun\c{c}\~{o}es e a verifique explicitamente (fazendo um c\'{a}lculo independente da teoria de Sturm-Liouville).
        \begin{solution}
            % TODO Escrever solu\c{c}\~{a}o.
        \end{solution}
    \end{parts}

    \question[E de 2010] Considere a equa\c{c}\~{a}o diferencial $y'' - y' + \lambda y = 0$ no intervalo $[0, 1]$.
    \begin{parts}
        \part Determine os autovalores e autofun\c{c}\~{o}es correspondentes para o problema de Sturm-Liouville associados \`{a} EDO acima sujeita \`{a}s condi\c{c}\~{o}es de controno $y(0) = y(1) = 0$.
        \begin{solution}
            % TODO Escrever solu\c{c}\~{a}o.
        \end{solution}

        \part Escreva a condi\c{c}\~{a}o de ortogonalidade entre essas autofun\c{c}\~{o}es.
        \begin{solution}
            % TODO Escrever solu\c{c}\~{a}o.
        \end{solution}
    \end{parts}

    \question[P2 de 2011, E de 2011] Encontre os autovalores e autofun\c{c}\~{o}es do seguinte problema de Sturm-Liouville:
    \begin{align*}
        \begin{cases}
            x(x y')' + \lambda y = 0, & 1 < x < \exp(2\pi) \\
            y'(1) = y'(\exp(2\pi)) = 0.
        \end{cases}
    \end{align*}
    Escreva a rela\c{c}\~{a}o de orotogonalidade satisfeita por essas autofun\c{c}\~{o}es.
    \begin{solution}
        Temos que $y = x^r$ de modo que
        \begin{align*}
            r (r - 1) + r + \lambda = 0
        \end{align*}
        e portanto $ r = \pm \sqrt{- \lambda}$.

        Para $\lambda < 0$ fazemos $\lambda = - k^2$, $k > 0$. Temos ent\~{a}o que $r = \pm k$ e portanto
        \begin{align*}
            y(x) &= A x^k + B x^{-k}, \\
            y'(x) &= A k x^{k - 1} - B k x^{k - 1}.
        \end{align*}
        Quando $x = 1$,
        \begin{align*}
            y'(1) = 0 \Rightarrow A + B = 0.
        \end{align*}
        Quando $x = \exp(2\pi)$,
        \begin{align*}
            y'(\exp(2\pi)) = 0 \Rightarrow A = 0.
        \end{align*}
        Pelas equa\c{c}\~{o}es acima temos que $A = B = 0$ e assim concluimos que para $\lambda < 0$ temos apenas a solu\c{c}\~{a}o trivial.

        Para $\lambda = 0$ temos
        \begin{align*}
            y(x) &= A_1 + B_1 \ln(x), \\
            y'(x) & = B_1 x^{-1}.
        \end{align*}
        Quando $x = 1$,
        \begin{align*}
            y'(1) = 0 \Rightarrow B_1 = 0.
        \end{align*}
        Quando $x = \exp(2\pi)$,
        \begin{align*}
            y'(\exp(2\pi)) = 0.
        \end{align*}
        Pelas equa\c{c}\~{o}es acima temos que $y = A_1$ para $\lambda = 0$.

        Para $\lambda > 0$ fazemos $\lambda = k^2$, $k > 0$. Temos ent\~{a}o que $r = \pm i k$ e portanto
        \begin{align*}
            y(x) &= A_2 x^{i k} + B_2 x^{-i k} \\
            &= A_2 \exp(i k \ln(x)) + B_2 \exp(-i k \ln(x)) \\
            &= C_1 \cos(k \ln(x)) + C_2 \sin(k \ln(x)), \\
            y'(x) &= -k x^{-1} C_1 \sin(k \ln(x)) + k x^{-1} C_2 \cos(k \ln(x)).
        \end{align*}
        Quando $x = 1$,
        \begin{align*}
            y'(1) = 0 \Rightarrow C_2 = 0.
        \end{align*}
        Quando $x = \exp(2\pi)$,
        \begin{align*}
            y'(\exp(2\pi)) = 0 \Rightarrow \sin(k \ln(\exp(2\pi))) = \sin(k 2 \pi) = 0 \Rightarrow 2 k \pi = n \pi, n = 1, 2, \ldots.
        \end{align*}
        Pelas equa\c{c}\~{o}es acima temos que $y_n = \cos(n 2^{-1} \ln(x))$, $n = 1, 2, \ldots$, e $\lambda_n = n^2 / 4$.

        Associando os casos em que $\lambda = 0$ e $\lambda > 0$ temos que
        \begin{align*}
            \begin{cases}
                y_n = \cos\left( n 2^{-1} \ln(x) \right), \\
                \lambda_n = n^2 / 4, n = 0, 1, 2, \ldots.
            \end{cases}
        \end{align*}

        Quanto a rela\c{c}\~{a}o de ortogonalidade temos
        \begin{align*}
            \int_1^{\exp(2\pi)} \cos(n 2^{-1} \ln(x)) \cos(m 2^{-1} \ln(x)) x^{-1} \,\mathrm{d}x = N_n \delta_{mn}, m, n = 0, 1, 2, \ldots.
        \end{align*}
    \end{solution}
\end{questions}
\bibliographystyle{plain}
\bibliography{bibliography}
\end{document}
